% Informationen ------------------------------------------------------------
% 	Definition von globalen Parametern, die im gesamten Dokument verwendet
% 	werden können (z.B auf dem Deckblatt etc.).
% --------------------------------------------------------------------------
\newcommand{\titel}{Dokumentation}
\newcommand{\untertitel}{Praktikum Robotik}
\newcommand{\art}{}
\newcommand{\fachgebiet}{Informatik}
\newcommand{\autor}{Bj\"orn Borgmann, Norman Caspari, Matthias Larisch und Philip Weber}
\newcommand{\studienbereich}{Informatik}
\newcommand{\matrikelnr}{}
\newcommand{\erstgutachter}{}
\newcommand{\zweitgutachter}{}
\newcommand{\jahr}{2012}

% Eigene Befehle
\newcommand{\todo}[1]{\textbf{\textsc{\textcolor{red}{(TODO: #1)}}}}

% Autorennamen in small caps
\newcommand{\AutorZ}[1]{\textsc{#1}}
\newcommand{\Autor}[1]{\AutorZ{\citeauthor{#1}}}

% Befehle zur semantischen Auszeichnung von Text
\newcommand{\NeuerBegriff}[1]{\textbf{#1}}
\newcommand{\Fachbegriff}[1]{\textit{#1}}
\newcommand{\Prozess}[1]{\textit{#1}}
\newcommand{\Webservice}[1]{\textit{#1}}
\newcommand{\Eingabe}[1]{\texttt{#1}}
\newcommand{\Code}[1]{\texttt{#1}}
\newcommand{\Datei}[1]{\texttt{#1}}
\newcommand{\Datentyp}[1]{\textsf{#1}}
\newcommand{\XMLElement}[1]{\textsf{#1}}

% Abkürzungen
\newcommand{\vgl}{Vgl.\ }
\newcommand{\ua}{\mbox{u.\,a.\ }}
\newcommand{\zB}{\mbox{z.\,B.\ }}
\newcommand{\bs}{$\backslash$}



% Dokumentenkopf -----------------------------------------------------------
% 	Diese Vorlage basiert auf "scrreprt" aus dem koma-script.
%		Die Option draft sollte beim fertigen Dokument ausgeschaltet werden.
% --------------------------------------------------------------------------
\documentclass[
	12pt,						% Schriftgr��e
	DIV10,
	german,					% f�r Umlaute, Silbentrennung etc.
	a4paper,					% Papierformat
	oneside,					% einseitiges Dokument
	titlepage,				% es wird eine Titelseite verwendet
	halfparskip,			% Abstand zwischen Abs�tzen (halbe Zeile)
	normalheadings,		% Gr��e der �berschriften verkleinern
	liststotoc,				% Verzeichnisse im Inhaltsverzeichnis auff�hren
	bibtotoc,				% Literaturverzeichnis im Inhaltsverzeichnis auff�hren
	idxtotoc,				% Index im Inhaltsverzeichnis auff�hren
	tablecaptionabove,	% Beschriftung von Tabellen oberhalb ausgeben
	final						% Status des Dokuments (final/draft)
]{scrreprt}
% Anpassung des Seitenlayouts ----------------------------------------------
% 	siehe Seitenstil.tex
% --------------------------------------------------------------------------
\usepackage[
	automark,			% Kapitelangaben in Kopfzeile automatisch erstellen
	headsepline,	% Trennlinie unter Kopfzeile
	ilines				% Trennlinie linksb�ndig ausrichten
]{scrpage2}


% Anpassung an Landessprache -----------------------------------------------
% 	Verwendet globale Option german siehe \documentclass
% --------------------------------------------------------------------------
\usepackage{babel}


% Umlaute ------------------------------------------------------------------
% 		Umlaute/Sonderzeichen wie ���� direkt im Quelltext verwenden (CodePage).
%		Erlaubt automatische Trennung von Worten mit Umlauten.
% --------------------------------------------------------------------------
\usepackage[utf8]{inputenc}
\usepackage[T1]{fontenc}
%\usepackage{ae} % "sch�neres" �
%\usepackage{textcomp} % Euro-Zeichen etc.

% Grafiken -----------------------------------------------------------------
% 		Einbinden von Grafiken [draft oder final]
% 		Option [draft] bindet Bilder nicht ein - auch globale Option
% --------------------------------------------------------------------------
%\usepackage[dvips,final]{graphicx}
\usepackage{graphicx}
\graphicspath{{Bilder/}} % Dort liegen die Bilder des Dokuments

% Zum einbinden von PDFs
\usepackage{pdfpages}

% Befehle aus AMSTeX f�r mathematische Symbole z.B. \boldsymbol \mathbb ----
\usepackage{amsmath,amsfonts}

% F�r Index-Ausgabe; \printindex -------------------------------------------
\usepackage{makeidx}

% Einfache Definition der Zeilenabst�nde und Seitenr�nder etc. -------------
\usepackage{setspace}
\usepackage{geometry}


% Symbolverzeichnis --------------------------------------------------------
% 	Symbolverzeichnisse bequem erstellen, beruht auf MakeIndex.
% 		makeindex.exe %Name%.nlo -s nomencl.ist -o %Name%.nls
% 	erzeugt dann das Verzeichnis. Dieser Befehl kann z.B. im TeXnicCenter
%		als Postprozessor eingetragen werden, damit er nicht st�ndig manuell
%		ausgef�hrt werden muss.
%		Die Definitionen sind ausgegliedert in die Datei Abkuerzungen.tex.
% --------------------------------------------------------------------------
\usepackage{nomencl}
  \let\abbrev\nomenclature
  \renewcommand{\nomname}{Abk�rzungsverzeichnis}
  \setlength{\nomlabelwidth}{.25\hsize}
  \renewcommand{\nomlabel}[1]{#1 \dotfill}
  \setlength{\nomitemsep}{-\parsep}


% Zum Umflie�en von Bildern -------------------------------------------------
%\usepackage{floatflt}


% Zum Einbinden von Programmcode --------------------------------------------
\usepackage{listings}
\usepackage{xcolor} 
\definecolor{hellgelb}{rgb}{1,1,0.9}
\definecolor{colKeys}{rgb}{0,0,1}
\definecolor{colIdentifier}{rgb}{0,0,0}
\definecolor{colComments}{rgb}{1,0,0}
\definecolor{colString}{rgb}{0,0.5,0}
\lstset{%
    float=hbp,%
    basicstyle=\texttt\small, %
    identifierstyle=\color{colIdentifier}, %
    keywordstyle=\color{colKeys}, %
    stringstyle=\color{colString}, %
    commentstyle=\color{colComments}, %
    columns=flexible, %
    tabsize=2, %
    frame=single, %
    extendedchars=true, %
    showspaces=false, %
    showstringspaces=false, %
    numbers=left, %
    numberstyle=\tiny, %
    breaklines=true, %
    backgroundcolor=\color{hellgelb}, %
    breakautoindent=true, %
%    captionpos=b%
}

% Lange URLs umbrechen etc. -------------------------------------------------
\usepackage{url}


% Wichtig f�r korrekte Zitierweise ------------------------------------------
\usepackage[square]{natbib}
% Quellenangaben in eckige Klammern fassen ----------------------------------
\bibpunct{[}{]}{;}{a}{}{,~}


% PDF-Optionen --------------------------------------------------------------
\usepackage[
bookmarks,
bookmarksopen=true,
pdftitle={\titel},
pdfauthor={\autor},
pdfcreator={\autor},
pdfsubject={\titel},
pdfkeywords={\titel},
colorlinks=true,
linkcolor=black, % einfache interne Verkn�pfungen
anchorcolor=black,% Ankertext
citecolor=blue, % Verweise auf Literaturverzeichniseintr�ge im Text
filecolor=magenta, % Verkn�pfungen, die lokale Dateien �ffnen
menucolor=black, % Acrobat-Men�punkte
urlcolor=cyan, 
% f�r die Druckversion k�nnen die Farben ausgeschaltet werden:
%linkcolor=black, % einfache interne Verkn�pfungen
%anchorcolor=black,% Ankertext
%citecolor=black, % Verweise auf Literaturverzeichniseintr�ge im Text
%filecolor=black, % Verkn�pfungen, die lokale Dateien �ffnen
%menucolor=black, % Acrobat-Men�punkte
%urlcolor=black, 
backref,
%pagebackref,
plainpages=false,% zur korrekten Erstellung der Bookmarks
pdfpagelabels,% zur korrekten Erstellung der Bookmarks
hypertexnames=false,% zur korrekten Erstellung der Bookmarks
linktocpage % Seitenzahlen anstatt Text im Inhaltsverzeichnis verlinken
]{hyperref}

% Zum fortlaufenden Durchnummerieren der Fu�noten ---------------------------
%\usepackage{chngcntr}


% f�r lange Tabellen
\usepackage{longtable}
\usepackage{array}
\usepackage{ragged2e}
\usepackage{lscape}
% f�r Tabellen mit variabler Spaltenbreite
%\usepackage{tabularx}

% Spaltendefinition rechtsb�ndig mit definierter Breite ---------------------
\newcolumntype{w}[1]{>{\raggedleft\hspace{0pt}}p{#1}}

% Formatierung von Listen �ndern
\usepackage{paralist}
% Standardeinstellungen:
% \setdefaultleftmargin{2.5em}{2.2em}{1.87em}{1.7em}{1em}{1em}

%bringt das H Kommando zur Bild Plazierung mit
%\usepackage{here} 

%\usepackage[utf8]{inputenc}

% Erstellung eines Index und Abk�rzungsverzeichnisses aktivieren -----------
\makeindex
%\makenomenclature


% Kopf- und Fu�zeilen, Seitenr�nder etc. -----------------------------------
% Zeilenabstand ------------------------------------------------------------
\onehalfspacing

% Seitenr�nder -------------------------------------------------------------
\geometry{paper=a4paper,left=35mm,right=35mm,top=30mm}

% Kopf- und Fu�zeilen ------------------------------------------------------
\pagestyle{scrheadings}

% Kopf- und Fu�zeile auch auf Kapitelanfangsseiten -------------------------
\renewcommand*{\chapterpagestyle}{scrheadings}

% Schriftform der Kopfzeile ------------------------------------------------
\renewcommand{\headfont}{\normalfont}

% Kopfzeile ----------------------------------------------------------------
\ihead{\large{\textsc{Seminararbeit}}\\	\small{Hierachisches Clustering interstellarer Objekte} \\[2ex] \textit{\headmark}}
\chead{}
\ohead{\includegraphics[scale=0.4]{UniOlLogoKlein.png}}
\setlength{\headheight}{21mm} % H�he der Kopfzeile
\setheadwidth[0pt]{textwithmarginpar} % Kopfzeile �ber den Text hinaus verbreitern
\setheadsepline[text]{0.4pt} % Trennlinie unter Kopfzeile

% Fu�zeile -----------------------------------------------------------------
\ifoot{}
\cfoot{}
\ofoot{\pagemark}

% erzeugt ein wenig mehr Platz hinter einem Punkt --------------------------
\frenchspacing 

% Schusterjungen und Hurenkinder vermeiden
\clubpenalty = 10000
\widowpenalty = 10000 
\displaywidowpenalty = 10000

% Quellcode-Ausgabe formatieren --------------------------------------------
\lstset{numbers=left, numberstyle=\tiny, numbersep=5pt, breaklines=true}
\lstset{emph={square}, emphstyle=\color{red}, emph={[2]root,base}, emphstyle={[2]\color{blue}}}

% Fu�noten fortlaufend durchnummerieren ------------------------------------
%\counterwithout{footnote}{chapter}

% kein page break bei chapter

\makeatletter
\renewcommand\chapter{\par%
    \global\@topnum\z@
    \@afterindentfalse
    \secdef\@chapter\@schapter}
\makeatother



% Eigene Definitionen f�r Silbentrennung
\hyphenation{Trenn-bar-es}
\hyphenation{J\"orgen}


% Das eigentliche Dokument -------------------------------------------------
%		Der eigentliche Inhalt des Dokuments beginnt hier. Die einzelnen Seiten
%		und Kapitel werden in eigene Dateien ausgelagert und hier nur inkludiert.
% --------------------------------------------------------------------------
\begin{document}

% auch subsubsection nummerieren
\setcounter{secnumdepth}{3}
\setcounter{tocdepth}{3}

% keine Kopf-/Fu�zeilen bei Deckblatt und Abstract
\ofoot{}
\thispagestyle{plain}
\begin{titlepage}

\begin{center}

\huge{\textsc{\textbf{\titel}}}\\[1.5ex]
\LARGE{\textbf{\untertitel}}\\[4ex]
\LARGE{\textbf{\art}}\\[6ex]%[1ex]
%\Large{im Fachgebiet \fachgebiet}\\[5ex]

\includegraphics[scale=0.25]{UniOlLogo.png}\\[2ex]

\normalsize
\begin{longtable}{rp{2.0ex}p{0.8\linewidth}}\\
%\begin{tabularx}{0.9\linewidth}{rp{2.0ex}X}\\
 %vorgelegt von:	& & \autor\\[1.2ex]
 Erstellt von:	& & \autor\\[1.2ex] %NEU
 Seminar: & & Data Mining in Astronomy\\[1.2ex]
 Studienbereich: & &\studienbereich\\[1.2ex]%[1.2ex]
 %Matrikelnummer: && \matrikelnr \\[1.2ex]
 %Erstgutachter:    &     & \erstgutachter \\[1.2ex]
 Professor:         && \erstgutachter\\[1.2ex]%[1.2ex] %NEU
 Tutor:     &    & \zweitgutachter \\[1.1ex]
%\end{tabularx}
\end{longtable}

%\copyright\ \jahr\\[1.5ex]

\end{center}

\singlespacing
\small
\noindent Dieses Werk einschließlich seiner Teile ist \textbf{urheberrechtlich geschützt}. Jede Verwertung außerhalb der engen Grenzen des Urheberrechtgesetzes ist ohne Zustimmung des Autors unzulässig und strafbar. Das gilt insbesondere für Vervielfältigungen, Übersetzungen, Mikroverfilmungen sowie die Einspeicherung und Verarbeitung in elektronischen Systemen.

\end{titlepage}

\ofoot{\pagemark}

% Seitennummerierung -------------------------------------------------------
%		Vor dem Hauptteil werden die Seiten in gro�en r�mischen Ziffern 
%		nummeriert...
% --------------------------------------------------------------------------
\pagenumbering{Roman}
\tableofcontents			% Inhaltsverzeichnis

% ...danach in normalen arabischen Ziffern ---------------------------------
\clearpage
\pagenumbering{arabic}

% Inhalt -------------------------------------------------------------------
%		Hier k�nnen jetzt die einzelnen Kapitel inkludiert werden. Sie m�ssen
%		in den entsprechenden .TEX-Dateien vorliegen. Die Dateinamen k�nnen
% 		nat�rlich angepasst werden.
% --------------------------------------------------------------------------
\chapter{Einleitung}

Diese Dokumentation beschreibt wie die, im Rahmen des Robotik Praktikums des SS 2012 entstandenen ROS Komponenten in eine ROS Installation integriert und verwendet werden können.

\paragraph{Aufgabenstellung}

\section{Umsetzung}
Die Aufgabenstellung gliedert das Projekt in mehrere Teilaufgaben, welche getrennt voneinander umgesetzt werden.

\subsection{Laserscanner}
Der Laserscanner LZR-U901 benötigt einen Treiber, welcher von der Universität Bremen erhältlich ist.

Dieser Treiber liest alle vier Scan-Ebenen des Laserscanners ein und gibt diese sowohl als PointCloud2 als auch die erste Ebene als LaserScan in ROS aus.

Um bei der Navigation Hindernissen besser ausweichen zu können, ist es sinnvoll, auf mehr als eine Ebene zurückzugreifen. Hierzu wird der Treiber so erweitert, dass zusätzlich ein zweites Topic \emph{merged_scan} veröffentlicht wird, welches eine nach bestimmten Regeln erstellte Ebene darstellt.

Da das Ziel der Zusammenlegung der Ebenen auf Hinderniserkennung festgelegt ist, wird immer der nächstliegende Punkt aus allen vier Ebenen in die zusammengelegte Ebene eingefügt.

Die Datenrate des Laserscanners beträgt für alle 4 Ebenen 3,75 Bilder pro Sekunde bei einer Winkelauflösung von 0.3516°. Die Ebenen liegen im Winkel von 2° zueinander.



\subsection{Kinect}
Während der Entwicklung stellten sich schnell Probleme mit der Rechenleistung des Turtlebot-Netbooks heraus. Als einfache Gegenmaßnahme wird die Kinect mit verringerter Auflösung- und Framerate betrieben. Das Tiefenbild wird mit 160x120 Pixeln und das RGB Bild mit 640x480 Pixeln bei jeweils 15 Bildern pro Sekunde verarbeitet.

\subsection{Navigation}
Zur Navigation wird die ROS node move_base verwendet. Die Einstellungen aus dem mitgelieferten turtlebot_navigation werden übernommen. Änderungen dienen hauptsächlich dazu, den Laserscanner sowie die Kinect als Sensoren zur Hinderniserkennung einzusetzen. Zudem werden sensorspezifische Konfigurationsdaten wie z.B. Reichweite entsprechend gesetzt.

Damit die Kartenerstellung besser funktioniert, werden zudem einige Bewegungsparameter so verändert, dass langsamere und weniger ruckartige Bewegungen erfolgen.


% Literaturverzeichnis -----------------------------------------------------
%		Das Literaturverzeichnis wird aus der Datenbank Bibliographie.bib 
% 		erstellt. Die genaue Verwendung von bibtex wird hier jedoch nicht erkl�rt.
%		Link: http://de.wikipedia.org/wiki/BibTeX
% --------------------------------------------------------------------------
\bibliography{Bibliographie}
\bibliographystyle{natdin}		% DIN-Stil des Literaturverzeichnisses

%\include{Inhalt/Erklaerung}	% Selbst�ndigkeitserkl�rung (wird hier nicht gebraucht)

% Anhang -------------------------------------------------------------------
%		Die Inhalte des Anhangs werden analog zu den Kapiteln inkludiert.
%		Dies geschieht in der Datei Anhang.tex
% --------------------------------------------------------------------------
\begin{appendix}
	\clearpage
	\pagenumbering{roman}
	%\input{Inhalt/Anhang} %% F�r Anhang Kommentar hier entfernen
\end{appendix}


% Index --------------------------------------------------------------------
%		Zum Erstellen eines Index, die folgende Zeile auskommentieren.
% --------------------------------------------------------------------------
%\printindex		% Index hier einf�gen

\end{document}
